
\section{Objetivo}
\label{inicio_cap4}

El presente capítulo establece los lineamientos para la gestión de las competencias, el entrenamiento, la evaluación y el mantenimiento de la idoneidad del personal que participa en las operaciones UAS desarrolladas por Skyfilms S.A.S.

Su propósito es asegurar que todo el personal operativo posea y mantenga los conocimientos, habilidades y experiencia requeridos para desempeñar sus funciones de forma segura, eficiente y conforme a los requisitos establecidos en el Reglamento Aeronáutico de Colombia (RAC 100), las Especificaciones Operacionales (OpSpecs) y los procedimientos internos de la organización.

Asimismo, este capítulo define los principios generales para la planificación del entrenamiento inicial, periódico y extraordinario, la evaluación continua del desempeño y la aplicación de los factores humanos como elemento fundamental para el fortalecimiento de la seguridad operacional.

\section{Competencias del Personal}

\subsection{Gerente General}


El Gerente General es el responsable de garantizar que Skyfilms S.A.S. disponga del personal, los recursos y la estructura organizacional necesarios para desarrollar operaciones UAS de forma segura y conforme a los requisitos establecidos por la Aerocivil de Colombia.

Como parte de sus competencias, deberá comprender las responsabilidades inherentes al explotador UAS, promover una cultura de seguridad operacional, aprobar los recursos destinados al entrenamiento del personal y asegurar el cumplimiento de las políticas establecidas en el presente Manual de Operaciones.

El Gerente General deberá mantener un conocimiento general de la normativa aplicable, del Sistema de Gestión de la Seguridad Operacional (SMS), del Manual de Control de Mantenimiento (MCM) y de las limitaciones operacionales aprobadas para la organización.



\subsection{Jefe de Pilotos}

El Jefe de Pilotos será responsable de supervisar la competencia técnica del personal operativo, coordinar los programas de entrenamiento y verificar que las operaciones UAS se desarrollen conforme al presente Manual de Operaciones y a las Especificaciones Operacionales aprobadas.

Deberá demostrar experiencia operacional suficiente en operaciones UAS, conocer la reglamentación aeronáutica aplicable, mantener vigente su competencia como piloto remoto y promover la estandarización de los procedimientos operacionales dentro de la organización.

Asimismo, será responsable de programar las evaluaciones periódicas del personal operativo, verificar la aplicación de los procedimientos establecidos y autorizar la participación del personal en las operaciones conforme a sus competencias demostradas.

\subsection{Pilotos Remotos}

Los Pilotos Remotos deberán demostrar los conocimientos, habilidades y experiencia necesarios para operar las aeronaves UAS autorizadas por Skyfilms S.A.S. de forma segura, eficiente y conforme a la reglamentación vigente.

Antes de participar en una operación deberán haber completado satisfactoriamente el entrenamiento inicial establecido por la organización, conocer el contenido del presente Manual de Operaciones, el Manual de Control de Mantenimiento, las Especificaciones Operacionales y los procedimientos aplicables al tipo de misión asignada.

Los Pilotos Remotos serán responsables de mantener vigente su competencia operacional mediante la participación en los programas de entrenamiento, las evaluaciones periódicas y los procesos de reentrenamiento establecidos por la organización.
    
\subsection{Observadores Visuales}

Los Observadores Visuales deberán recibir entrenamiento específico para apoyar la ejecución segura de las operaciones UAS cuando su participación sea requerida por las características de la misión o por las condiciones operacionales.

Su competencia comprenderá la vigilancia permanente del espacio aéreo y del entorno operacional, la identificación de peligros, la comunicación efectiva con el Piloto Remoto y la aplicación de los procedimientos establecidos para la prevención de conflictos con personas, aeronaves y obstáculos.

Los Observadores Visuales deberán conocer los procedimientos de emergencia aplicables y participar en los entrenamientos y evaluaciones definidos por Skyfilms S.A.S.

\subsection{Responsable de Mantenimiento}

El Responsable de Mantenimiento deberá poseer los conocimientos técnicos necesarios para garantizar que las aeronaves UAS y sus equipos asociados permanezcan en condiciones de aeronavegabilidad y disponibilidad operacional.

Será responsable de supervisar el cumplimiento del Manual de Control de Mantenimiento (MCM), coordinar las inspecciones programadas, verificar la correcta ejecución de las actividades de mantenimiento preventivo y correctivo, controlar la trazabilidad de los componentes críticos y administrar los registros técnicos de cada aeronave.

Asimismo, deberá mantener actualizados sus conocimientos sobre las aeronaves operadas por Skyfilms S.A.S., sus sistemas, firmware, baterías, enlaces de mando y control y demás componentes que puedan afectar la seguridad operacional.


\subsection{Personal Operativo Vinculado}
A continuación, se relaciona el personal operativo autorizado por Skyfilms S.A.S. Todo Piloto UAS relacionado en esta tabla cuenta con las competencias, certificaciones y entrenamiento continuo exigidos por la Aerocivil.

\begin{tcolorbox}[title=Nota Operacional,colback=gray!5,colframe=black]
De conformidad con el RAC 100, todo Piloto UAS titular de un Certificado de Idoneidad (CIPU) vigente, se encuentra plenamente facultado y cualificado para ejercer el rol de Observador Visual (Observador UA) en cualquier misión que así lo requiera.

\end{tcolorbox}

\begin{table}[H]
    \centering
    \renewcommand{\arraystretch}{1.2}
    \begin{tabular}{|p{4cm}|p{2.5cm}|p{2.5cm}|p{2cm}|p{3.5cm}|}
        \hline
        \rowcolor[gray]{0.9} \textbf{Nombre} & \textbf{Cédula} & \textbf{Cargo} & \textbf{CIPU} & \textbf{ Rol / Adiciones} \\ \hline
        Andrés Felipe Pérez Sierra & 1108206142 & Jefe de Pilotos & 100-24-2431 & Piloto y Observador / Nocturno \\ \hline
        Cesar León Castrillón Castrillón & 71362721 & Piloto UAS & 100-24-2057 &Piloto y Observador / Nocturno \\ \hline
        % Agrega más filas si tienes otros pilotos
    \end{tabular}
\end{table}

\textit{Las copias de los certificados de idoneidad (CIPU) y certificados médicos del personal se adjuntan en los Anexos de este manual y reposan en el Sistema de Gestión Documental de la empresa.}



\section{Programa de Entrenamiento}

Skyfilms S.A.S. establecerá y mantendrá un Programa de Entrenamiento destinado a garantizar que todo el personal involucrado en las operaciones UAS adquiera, mantenga y actualice las competencias necesarias para el desempeño seguro de sus funciones.

El Programa de Entrenamiento comprenderá actividades de formación inicial, entrenamiento periódico, entrenamiento asociado a cambios operacionales y entrenamiento extraordinario cuando las circunstancias operacionales lo requieran.

Los contenidos, metodologías, evaluaciones y registros del entrenamiento se desarrollarán en el Programa de Entrenamiento de la organización, el cual constituye un documento controlado complementario al presente Manual de Operaciones.

\subsection{Entrenamiento Inicial}

Todo el personal operativo deberá completar satisfactoriamente un entrenamiento inicial antes de asumir funciones dentro de las operaciones UAS de Skyfilms S.A.S.

El entrenamiento inicial tendrá como finalidad proporcionar los conocimientos teóricos y prácticos necesarios para el desempeño seguro de las funciones asignadas, incluyendo el conocimiento de la organización, la normativa aplicable, los procedimientos operacionales, las limitaciones de las aeronaves, los procedimientos de emergencia y los principios básicos de seguridad operacional.

Ningún miembro del personal podrá participar en operaciones UAS sin haber completado el entrenamiento inicial correspondiente y haber sido declarado competente por el Jefe de Pilotos o por el instructor designado.

\subsection{Entrenamiento Periódico}

Skyfilms S.A.S. desarrollará actividades de entrenamiento periódico con el propósito de mantener las competencias operacionales del personal, actualizar conocimientos y reforzar la correcta aplicación de los procedimientos establecidos por la organización.

La periodicidad, alcance y contenido del entrenamiento serán definidos en el Programa de Entrenamiento, considerando la experiencia del personal, los cambios normativos, las lecciones aprendidas, los resultados de auditorías internas y los eventos reportados dentro del Sistema de Gestión de la Seguridad Operacional (SMS).

\subsection{Entrenamiento por Cambios Operacionales}

Siempre que se incorporen nuevas aeronaves, equipos, tecnologías, procedimientos, escenarios operacionales o modificaciones regulatorias que puedan afectar la seguridad operacional, Skyfilms S.A.S. desarrollará el entrenamiento correspondiente antes de autorizar su utilización en operaciones.

Este entrenamiento tendrá como propósito garantizar que el personal comprenda los cambios implementados y demuestre la competencia necesaria para su aplicación segura.

\subsection{Entrenamiento Extraordinario}

Skyfilms S.A.S. podrá programar entrenamientos extraordinarios cuando se identifiquen deficiencias operacionales, incumplimientos de procedimientos, cambios significativos en la reglamentación, incorporación de nuevos riesgos operacionales o como resultado de investigaciones de incidentes, accidentes o hallazgos derivados de auditorías.

El alcance y contenido del entrenamiento extraordinario serán definidos por el Jefe de Pilotos en coordinación con el Responsable del Sistema de Gestión de la Seguridad Operacional (SMS), cuando aplique.

%%%%%%%%%%%%%%%%%%%%%%%%%%%%%%%%
\section{Desarrollo del Programa de Entrenamiento}


La formación del personal de Skyfilms S.A.S. se desarrollará mediante un enfoque basado en competencias, orientado a garantizar que cada integrante de la organización adquiera y mantenga los conocimientos, habilidades y actitudes necesarias para el desempeño seguro de sus funciones.

El Programa de Entrenamiento combinará actividades teóricas y prácticas, adaptadas a las responsabilidades de cada cargo y al tipo de operación desarrollada por la organización. Los contenidos y la intensidad del entrenamiento serán definidos de acuerdo con los riesgos operacionales, la experiencia del personal, las características de las aeronaves utilizadas y las modificaciones que se introduzcan en la reglamentación o en los procedimientos internos.

\subsection{Contenido del Entrenamiento}

El Programa de Entrenamiento incluirá, según las funciones desempeñadas por cada integrante del personal, contenidos relacionados con la normativa aeronáutica vigente, políticas y procedimientos internos, planificación y ejecución de operaciones UAS, factores humanos, gestión de la seguridad operacional, procedimientos de emergencia, mantenimiento básico de las aeronaves, uso de los equipos asociados y demás materias necesarias para garantizar el cumplimiento de las funciones asignadas.

El contenido específico de cada curso, su duración y los criterios de aprobación serán establecidos en el Programa de Entrenamiento de Skyfilms S.A.S.





\subsection{Matriz de Capacitación}

La siguiente matriz relaciona los cursos del Programa de Entrenamiento con cada cargo. Se marca \textbf{I} (inicial, antes de operar) y \textbf{P} (periódico, cada 4 meses), y \textbf{--} cuando no aplica.

\begin{table}[H]
    \centering
    \renewcommand{\arraystretch}{1.2}
    \begin{tabular}{|p{6cm}|c|c|c|c|}
        \hline
        \rowcolor[gray]{0.9} \textbf{Curso} & \textbf{Jefe de Pilotos} & \textbf{Piloto Remoto} & \textbf{Observador} & \textbf{Mantenim.} \\ \hline
        Manual de Operaciones, OpSpecs y RAC 100 & I / P & I / P & I / P & I \\ \hline
        SMS y reporte de peligros (F-006, F-012) & I / P & I / P & I / P & I / P \\ \hline
        Planificación y evaluación del riesgo (F-009, F-011) & I / P & I / P & I & -- \\ \hline
        Operación del Mavic 4 Pro & I / P & I / P & -- & I \\ \hline
        Operación del Pavo20 Pro (FPV) & I / P & I / P & -- & I \\ \hline
        Enlace C2, pérdida de enlace y failsafe & I / P & I / P & I & I \\ \hline
        Vuelo nocturno: luces, dronpuerto y desorientación & I / P & I / P & I / P & -- \\ \hline
        Funciones del observador y comunicaciones & I / P & I / P & I / P & -- \\ \hline
        Emergencias: fuego de batería, aterrizaje forzoso, accidente (F-013) & I / P & I / P & I / P & I / P \\ \hline
        Baterías LiPo: carga, almacenamiento y transporte & I / P & I / P & I & I / P \\ \hline
        Factores humanos y fatiga & I / P & I / P & I / P & I \\ \hline
        Mantenimiento básico y registros (F-005, F-007) & I & I & -- & I / P \\ \hline
        Diligenciamiento de los formatos F-001 a F-017 & I / P & I / P & I & I / P \\ \hline
        Repaso de manuales de usuario de la aeronave y de los equipos & I / P & I / P & I & I / P \\ \hline
    \end{tabular}
\end{table}

\subsubsection*{Entrenamiento específico por aeronave y por tipo de vuelo}

\begin{itemize}
\item \textbf{Mavic 4 Pro:} manejo del control DJI RC Pro 2 y de DJI Fly, modos de vuelo, RTH, límites y alertas, antes del primer vuelo en operaciones.
\item \textbf{Pavo20 Pro (FPV):} el piloto remoto debe completar vuelos de práctica en modo de asistencia y en modo acro, configuración y prueba de failsafe, uso de las gafas y comunicación con el observador. El primer uso en operación requiere una evaluación práctica del Jefe de Pilotos.
\item \textbf{Vuelo nocturno:} todo piloto que opere de noche debe contar con la adición nocturna en su certificado de idoneidad (RAC 100.630) y con entrenamiento en el procedimiento de vuelo nocturno del Manual: iluminación de la aeronave y del dronpuerto, referencias visuales, desorientación y pérdida del enlace.
\item \textbf{Cambio de equipo:} una nueva aeronave, firmware, módulo de radio o configuración de failsafe requiere entrenamiento antes de operar (Entrenamiento por Cambios Operacionales).
\end{itemize}

\subsection{Métodos de Instrucción}

Skyfilms S.A.S. podrá emplear diferentes métodos de instrucción para el desarrollo de las actividades de entrenamiento, procurando que estos respondan a las necesidades operacionales de la organización y faciliten el desarrollo progresivo de las competencias del personal.

Los métodos de instrucción podrán comprender, entre otros:

\begin{itemize}
    \item Clases teóricas presenciales o virtuales.
    \item Autoestudio dirigido mediante documentación controlada.
    \item Talleres y ejercicios prácticos.
    \item Simulación de escenarios operacionales.
    \item Prácticas de vuelo supervisadas.
    \item Análisis de casos e incidentes operacionales.
    \item Entrenamiento en el puesto de trabajo (\textit{On the Job Training}).
\end{itemize}

La selección del método de instrucción dependerá de los objetivos del entrenamiento, el nivel de experiencia del personal y la complejidad de las actividades a desarrollar.




\subsection{Registro y Control del Entrenamiento}

Skyfilms S.A.S. mantendrá registros actualizados de todas las actividades de entrenamiento desarrolladas por su personal, con el propósito de demostrar el cumplimiento de los requisitos de competencia establecidos por la organización y facilitar el seguimiento de la vigencia de las habilitaciones internas.

Los registros incluirán, como mínimo, la identificación del participante, la actividad desarrollada, la fecha de realización, el instructor o evaluador responsable, los resultados obtenidos y las observaciones que resulten pertinentes.

El Jefe de Pilotos custodia y valida los registros de entrenamiento: verifica que cada F-016 y F-017 esté completo y firmado, que la vigencia de cada persona esté al día y que el cronograma se cumpla. Nadie opera si tiene el entrenamiento periódico vencido.

Cada capacitación se registra en el Registro de Asistencia (F-016) y se cierra con el formato de Evaluación y Cierre (F-017). El programa anual se publica con las fechas de cada ciclo de 4 meses; si una fecha cambia, el Jefe de Pilotos actualiza el cronograma y avisa al personal.

La administración, conservación y actualización de estos registros se realizará conforme al procedimiento de Control Documental establecido por la organización.

%%%%%%%%%%%%%%%%%%%%%%%%%%%%%%%%%%%%%%%%%%%%%%%%%%%%%%%%%%%%%%%%%%%%%%%%%%%%%%%%

\section{Evaluación y Mantenimiento de Competencias}

\subsection{Evaluaciones Teóricas}

Las evaluaciones teóricas tendrán como finalidad verificar el conocimiento del personal respecto a la normativa aeronáutica aplicable, el presente Manual de Operaciones, las Especificaciones Operacionales (OpSpecs), el Manual de Control de Mantenimiento (MCM), el Sistema de Gestión de la Seguridad Operacional (SMS), los procedimientos internos y demás documentos aplicables a sus funciones.

La metodología de evaluación será definida en el Programa de Entrenamiento y podrá comprender exámenes escritos, cuestionarios, entrevistas técnicas u otros mecanismos que permitan verificar el nivel de conocimiento requerido.

El personal deberá obtener los resultados mínimos establecidos por la organización para ser considerado competente.



\subsection{Evaluaciones Prácticas}

Las evaluaciones prácticas tendrán como propósito verificar la capacidad del personal para aplicar correctamente los procedimientos operacionales durante la planificación, ejecución y finalización de las operaciones UAS.

Las evaluaciones podrán incluir ejercicios de vuelo, simulaciones, resolución de escenarios operacionales, aplicación de procedimientos normales y de emergencia, así como cualquier otra actividad que permita evidenciar la competencia operacional del evaluado.

Las evaluaciones serán realizadas por el Jefe de Pilotos o por el personal designado por la organización para tal fin.

\subsection{Seguimiento del Desempeño Operacional}

Skyfilms S.A.S. realizará seguimiento continuo al desempeño operacional de su personal mediante la observación de las operaciones, el análisis de reportes operacionales, los resultados de auditorías internas, los registros de entrenamiento y la información proveniente del Sistema de Gestión de la Seguridad Operacional (SMS).

Cuando se identifiquen desviaciones respecto a los procedimientos establecidos, disminución en el nivel de competencia o cualquier condición que pueda afectar la seguridad operacional, el Jefe de Pilotos determinará las acciones correctivas correspondientes, las cuales podrán incluir entrenamiento adicional, reentrenamiento, supervisión operacional o suspensión temporal de la autorización para participar en operaciones.

%%%%%%%%%%%%%%%%%%%%%%%%%%%%%%%%%%%%%%%%%%%%%%%%%%%%%%%%%%%%%%%%%%%%%%%%%%%%%%%%
\section{Reentrenamiento}


Skyfilms S.A.S. establecerá un programa de reentrenamiento con el propósito de restablecer o reforzar las competencias del personal cuando se identifiquen condiciones que puedan afectar el desempeño seguro de las operaciones UAS.

El reentrenamiento tendrá como objetivo corregir deficiencias detectadas durante las evaluaciones, actualizar conocimientos, fortalecer habilidades operacionales y asegurar el cumplimiento permanente de los estándares establecidos por la organización.

El alcance, contenido y duración del reentrenamiento serán definidos por el Jefe de Pilotos de acuerdo con las necesidades identificadas y los resultados de las evaluaciones realizadas.



\subsection{Criterios para Reentrenamiento}

El personal podrá ser programado para un proceso de reentrenamiento cuando se presente cualquiera de las siguientes situaciones:

\begin{itemize}

\item Resultados insatisfactorios en las evaluaciones teóricas o prácticas.

\item Incumplimiento de los procedimientos operacionales establecidos por la organización.

\item Participación en incidentes o eventos operacionales que evidencien deficiencias en el desempeño.

\item Incorporación de nuevas aeronaves, tecnologías o procedimientos que requieran entrenamiento adicional.

\item Cambios en la reglamentación aeronáutica aplicable.

\item Ausencia prolongada de actividades operacionales que pueda afectar el mantenimiento de las competencias.

\item Cualquier otra condición identificada por el Jefe de Pilotos que justifique la necesidad de reforzar las competencias del personal.

\end{itemize}



\subsection{Suspensión Temporal de Competencias}

Cuando se determine que un integrante del personal no demuestra las competencias necesarias para desempeñar sus funciones de manera segura, Skyfilms S.A.S. podrá suspender temporalmente su autorización para participar en operaciones UAS.

La suspensión podrá aplicarse como resultado de evaluaciones insatisfactorias, incumplimientos graves de los procedimientos operacionales, limitaciones médicas conocidas, incumplimiento de los requisitos de entrenamiento o cualquier otra condición que represente un riesgo para la seguridad operacional.

Durante el periodo de suspensión el personal no podrá ejercer las funciones para las cuales fue suspendido hasta cumplir satisfactoriamente el proceso de reentrenamiento y las evaluaciones correspondientes.



\subsection{Reincorporación a las Operaciones}

La reincorporación del personal a las actividades operacionales estará condicionada al cumplimiento satisfactorio del proceso de reentrenamiento y a la demostración de las competencias requeridas para el desempeño seguro de sus funciones.

El Jefe de Pilotos verificará el cumplimiento de los requisitos establecidos y autorizará nuevamente la participación del personal en operaciones cuando considere que las competencias han sido restablecidas.

La autorización de reincorporación deberá quedar registrada conforme al procedimiento de control documental de la organización.


\begin{tcolorbox}[
title=Nota Operacional,
colback=gray!5,
colframe=black]

El reentrenamiento constituye una medida preventiva orientada al fortalecimiento de las competencias del personal y no deberá interpretarse como una medida disciplinaria. Su propósito es garantizar el mantenimiento de los estándares de seguridad operacional establecidos por Skyfilms S.A.S.

\end{tcolorbox}

%%%%%%%%%%%%%%%%%%%%%%%%%%%%%%%%%%%%%%%%%%%%%%%%%%%%%%%%%%%%%%%%%%%%%%%%%%%%%%%%

\section{Análisis de Factores Humanos}

Skyfilms S.A.S. reconoce que el desempeño humano constituye un elemento determinante para la seguridad operacional. En consecuencia, promoverá una cultura organizacional orientada a la prevención de errores, la toma de decisiones responsable, la comunicación efectiva y la mejora continua de las competencias del personal.

La gestión de los factores humanos se integrará a los procesos de entrenamiento, supervisión operacional y Sistema de Gestión de la Seguridad Operacional (SMS), con el propósito de reducir la probabilidad de errores y fortalecer la capacidad del personal para identificar, comunicar y gestionar los riesgos presentes durante las operaciones UAS.

\subsection{Cultura de Seguridad Operacional}

Skyfilms S.A.S. promoverá una cultura de seguridad operacional basada en el compromiso de todo el personal con el cumplimiento de la reglamentación aeronáutica, los procedimientos internos y las buenas prácticas operacionales.

Todo el personal tendrá la responsabilidad de actuar de manera segura, comunicar oportunamente los peligros identificados y contribuir al fortalecimiento continuo del desempeño operacional de la organización.

La seguridad operacional prevalecerá sobre cualquier consideración comercial, operacional o administrativa que pueda comprometer la integridad de las personas, las aeronaves o los bienes involucrados en las operaciones.

\subsection{Gestión del Error Humano}

Skyfilms S.A.S. reconoce que el error humano constituye una condición inherente a las actividades desarrolladas por las personas. En consecuencia, la organización orientará sus esfuerzos hacia la prevención, identificación y gestión de los errores antes que a la asignación de responsabilidades individuales.

Los errores identificados durante las operaciones serán analizados con el propósito de determinar sus causas contribuyentes, establecer acciones correctivas y prevenir su recurrencia, promoviendo una cultura de aprendizaje y mejora continua.

El procedimiento ante un error humano es el siguiente:
\begin{enumerate}
\item \textbf{Identificar:} quien lo detecte lo informa de inmediato al piloto remoto o al Jefe de Pilotos y, si afectó la seguridad, diligencia el F-006.
\item \textbf{Controlar:} el piloto remoto contiene la situación (pausar, aterrizar o suspender) antes de continuar.
\item \textbf{Analizar:} el Gerente SMS, con el Jefe de Pilotos, determina las causas y lo registra en el F-012.
\item \textbf{Corregir:} se define una acción correctiva con responsable y fecha, por ejemplo reentrenamiento o cambio de procedimiento.
\item \textbf{Seguir:} el Gerente SMS verifica que la acción se cumplió y archiva el registro; no se sanciona al reportante de buena fe.
\end{enumerate}

\subsection{Conciencia Situacional}

El personal operativo deberá mantener permanentemente una adecuada conciencia situacional durante todas las fases de la operación, identificando oportunamente los factores que puedan afectar la seguridad operacional.

La conciencia situacional comprenderá el conocimiento del entorno operacional, las condiciones meteorológicas, el espacio aéreo, la ubicación de personas y obstáculos, el estado de la aeronave, las comunicaciones y cualquier otro elemento que pueda influir sobre la ejecución segura de la misión.

\subsection{Toma de Decisiones Aeronáuticas}

Las decisiones operacionales deberán fundamentarse en una adecuada evaluación de riesgos, el cumplimiento de los procedimientos establecidos y el criterio profesional del personal responsable de la operación.

Cuando las condiciones operacionales no permitan garantizar un nivel aceptable de seguridad, el Piloto Remoto tendrá la autoridad para modificar, retrasar, interrumpir o cancelar la misión, informando oportunamente al Jefe de Pilotos cuando corresponda.

\subsection{Comunicación y Trabajo en Equipo}

La comunicación entre los integrantes del equipo operacional deberá ser clara, oportuna y verificable, asegurando que toda la información relevante para la seguridad de la operación sea comprendida por las personas involucradas.

Durante las operaciones que requieran la participación de Observadores Visuales u otro personal de apoyo, se emplearán procedimientos de comunicación previamente establecidos y conocidos por todos los participantes.

\subsection{Fatiga y Aptitud para el Servicio}

Todo el personal deberá presentarse a sus actividades en condiciones físicas y mentales compatibles con la ejecución segura de las operaciones.

Ningún integrante del personal podrá participar en operaciones cuando se encuentre afectado por fatiga excesiva, enfermedad, consumo de alcohol, sustancias psicoactivas, medicamentos que alteren sus capacidades o cualquier otra condición que comprometa su desempeño.

Cada integrante del personal tendrá la obligación de informar oportunamente cualquier condición que pueda afectar su aptitud para el servicio.

Para la prevención de la fatiga se aplican, como mínimo, los siguientes límites (RAC 100.230 y 100.510):

\begin{itemize}
\item No se opera dentro de las 12 horas siguientes al consumo de alcohol, ni bajo el efecto de sustancias psicoactivas.
\item El piloto remoto no realiza operaciones de vuelo durante un período de servicio superior a 10 horas continuas dentro de un período de 24 horas, independientemente del tiempo de vuelo efectivo.
\item Entre dos jornadas de servicio, el piloto y el observador descansan al menos 8 horas continuas; después de una jornada de 10 horas no se programa otra operación en el mismo día.
\item El Jefe de Pilotos programa las operaciones con pausas y verifica el descanso con el F-014; cualquier integrante puede declararse no apto y lo informa al Jefe de Pilotos, sin consecuencias disciplinarias.
\item Antes de cada operación, el piloto diligencia el factor de descanso y carga de trabajo del F-011; si lo califica como crítico, no opera (NO GO).
\end{itemize}

\subsection{Métodos de evaluación y retroalimentación}

Los factores humanos se evalúan con observación en vuelo por el Jefe de Pilotos, el análisis de los F-006 y F-012, las evaluaciones del F-017 y las conversaciones de cierre (\textit{debriefing}) después de cada misión. Cada cuatro meses, el Gerente SMS entrega al personal los hallazgos y las lecciones aprendidas dentro del entrenamiento periódico.

\subsection{Mejora Continua}

Skyfilms S.A.S. fomentará la mejora continua mediante el análisis de las experiencias operacionales, los reportes voluntarios, las lecciones aprendidas, los resultados de auditorías y las recomendaciones derivadas del Sistema de Gestión de la Seguridad Operacional (SMS).

La información obtenida será utilizada para fortalecer los programas de entrenamiento, actualizar los procedimientos operacionales y mejorar continuamente el desempeño de la organización.


\begin{tcolorbox}[
title=Principio de Seguridad Operacional,
colback=gray!5,
colframe=black]

En Skyfilms S.A.S., toda persona tiene la autoridad y la responsabilidad de detener una operación cuando considere que las condiciones de seguridad no son adecuadas. Ningún objetivo operativo, comercial o administrativo tendrá prioridad sobre la protección de las personas, la seguridad operacional y el cumplimiento de la reglamentación aeronáutica.

\end{tcolorbox}
